\begin{tcolorbox}[colback=propbox, colframe=propborder, arc=2pt,
  left=6pt, right=6pt, top=4pt, bottom=4pt,
  title={\small\textbf{Definition (Parse Tree)}},
  fonttitle=\small\bfseries]
\label{def:parse-tree}
The \emph{parse tree}, or \emph{formation tree}, of a well-formed formula is a
tree representation of its recursive construction:
\begin{itemize}
  \item leaves are propositional variables;
  \item a negation node has one child;
  \item a binary-connective node has two children.
\end{itemize}
\end{tcolorbox}

\begin{remark*}[Interpretation]
The parse tree records the same structure described by unique readability. It
turns the recursive construction of a formula into a visible tree.
\end{remark*}

\begin{dependencies}
\begin{itemize}
  \item \hyperref[def:well-formed-formula]{Well-Formed Formula}
  \item \hyperref[def:logical-connective]{Logical Connective}
  \item \hyperref[thm:unique-readability-propositional-logic]{Unique Readability for Propositional Logic}
\end{itemize}
\end{dependencies}
